% Sección (iv) Experimental Setup. Se inserta en main.tex antes de la Conclusión:
%   \clearpage
%   % Sección (iv) Experimental Setup. Se inserta en main.tex antes de la Conclusión:
%   \clearpage
%   % Sección (iv) Experimental Setup. Se inserta en main.tex antes de la Conclusión:
%   \clearpage
%   % Sección (iv) Experimental Setup. Se inserta en main.tex antes de la Conclusión:
%   \clearpage
%   \input{experimental_setup}
% Usa solo paquetes ya cargados por main.tex y pucv_inf_2024.sty.
% Referencia eq:objetivo, definida en la Propuesta metodológica.
\section{Configuración experimental}
\label{sec:experimental}

Este diseño responde a la pregunta del proyecto: ¿en qué condiciones D-MAB
mejora Tabu Search frente a la selección uniforme? Por ello, las comparaciones
comparten representación, planificador, objetivo y presupuesto. La hipótesis
primaria plantea que, bajo congestión, Tabu con D-MAB obtiene menor $f$ medio que
Tabu uniforme. El equipo fijó hipótesis, semillas y pruebas en un protocolo previo
a los experimentos nuevos. Esta sección actualiza el plan de evaluación de la sección
anterior y prevalece ante cualquier diferencia. Además, describe el diseño y no informa
conclusiones.

\subsection{Conjuntos de datos}
\label{subsec:datos}
El estudio usa dos catálogos sintéticos, sin datos clínicos, y los analiza por
separado. El catálogo estándar proviene del repositorio oficial
\parencite{saicooh2026}. Contiene 12 instancias: tres réplicas de 15, 20, 25 y 30
cirugías. Sin embargo, este catálogo casi no tensiona la espera entre etapas. Solo
16 de 240 soluciones finales del estudio de carteras presentan espera positiva.
Por tanto, no permite evaluar la condición de bloqueo que plantea la pregunta.

Para cubrir ese vacío, el equipo generó un catálogo de congestión de 48 instancias,
doce por tamaño. El generador remuestrea pacientes de una instancia estándar y
aplica ruido lognormal con $\sigma=0{,}10$. Luego redondea los tiempos a 0,5 minutos
y redibuja la elegibilidad de personal. Cada instancia tiene semilla propia.

Una calibración evaluó 18 combinaciones de salas, cirujanos y espera máxima, sobre
cuatro instancias de desarrollo excluidas del catálogo. La regla pre-registrada eligió
reducir los cirujanos al 40~\% y mantener salas y espera máxima. Esta elección
sigue el mecanismo del modelo: con la misma sala en ambas etapas, la espera depende
del cirujano. De hecho, sin reducir cirujanos, ninguna combinación superó 10~\% de
soluciones con espera. En el catálogo final, 51,8~\% de las soluciones de Tabu uniforme
y D-MAB presenta espera positiva. Este dato verifica la presión sobre la restricción.

Ambos catálogos superan la validación de esquema, tiempos, recursos, elegibilidad y
digest canónico. Además, FIFO, SPT y LPT encuentran solución factible en las 48
instancias nuevas. Los manifiestos guardan el SHA-256 de cada archivo. Cada
experimento usa un rango exclusivo de semillas. Las semillas nuevas sobre el
catálogo estándar miden variación de la búsqueda; no miden generalización a
instancias no vistas.

\subsection{Líneas base y controles}
\label{subsec:controles}
La Tabla~\ref{tab:lineas-base} vincula cada control con la explicación alternativa
que descarta. Así, una diferencia observada puede atribuirse al aprendizaje y no a
otro componente del diseño.

\begin{table}[ht]
 \centering
 \caption{Controles y explicación alternativa que descarta cada uno.}
 \label{tab:lineas-base}
 \begin{tabular}{@{}p{0.33\textwidth}p{0.61\textwidth}@{}}
  \toprule
  Control (evaluaciones) & Explicación alternativa\\
  \midrule
  FIFO, SPT y LPT & Una regla constructiva basta. Integran el inicio de Tabu; superarlas no prueba el aporte de D-MAB.\\
  Tabu uniforme (3.030) & La mejora proviene de Tabu y no de aprender.\\
  Tabu prolongado (6.060) & La cartera mejora solo por consumir más evaluaciones.\\
  Dos Tabu independientes (6.060) & La cartera mejora por repetir la búsqueda.\\
  Tabu + UCB1 sin Page--Hinkley (3.030) & El efecto proviene de UCB1 y no del reinicio.\\
  \bottomrule
 \end{tabular}
\end{table}

\subsection{Métricas}
\label{subsec:metricas}
La métrica primaria es $f(S)$ de la ecuación~\eqref{eq:objetivo}, fijada antes de
evaluar. Las métricas secundarias son $C_{\max}$, espera total, espera máxima,
factibilidad, tiempo de pared y convergencia. La convergencia se mide con la
evaluación de la última mejora global. Estas métricas son descriptivas. Una baja de
$f$ no implica una baja de cada componente; por eso, se informan por separado.

\subsection{Implementación}
\label{subsec:implementacion}
El código usa Python~3.13.2, NumPy~2.2.3, SciPy~1.15.2 y PyYAML~6.0.3. Los
experimentos nuevos corren en un AMD Ryzen~7 8845HS, con 8 núcleos, 16 hilos y
13,8~GiB de RAM, sobre Windows~11. La Tabla~\ref{tab:parametros} fija los parámetros.
El tiempo se mide en un solo proceso. Una cartera tarda, en mediana, 4,12~s con 15
cirugías (rango 3,83--4,80) y 13,09~s con 30 (rango 12,65--13,53). Sin embargo, el
rendimiento varió entre sesiones: con 15 cirugías pasó de 1.470 a cerca de 2.240
evaluaciones por segundo. Por eso, el control principal de cómputo son las evaluaciones.

\begin{table}[ht]
 \centering
 \caption{Parámetros congelados de la búsqueda y del presupuesto.}
 \label{tab:parametros}
 \begin{tabular}{@{}p{0.16\textwidth}p{0.78\textwidth}@{}}
  \toprule
  Componente & Configuración\\
  \midrule
  Tabu & 30 soluciones iniciales (FIFO, SPT, LPT y 27 aleatorias); lotes de 15 candidatos; memoria de 7 firmas de solución; aspiración por mejora global.\\
  D-MAB & Brazos \texttt{swap}, \texttt{insert}, \texttt{room\_1}, \texttt{room\_2} y \texttt{room\_pair}; $C=0{,}1$; $\delta=0{,}01$; $\lambda=1$.\\
  Cartera & Dos caminos de 3.030 evaluaciones; reparto 50/50 fijo; conserva el menor $f$.\\
  Parada & Agotar el presupuesto de evaluaciones del camino.\\
  Conteo & Cuentan las 30 iniciales y cada candidato distinto del lote, aun infactible. No cuentan duplicados, reparación ni revalidación final.\\
  \bottomrule
 \end{tabular}
\end{table}

El archivo de 30 soluciones del código no altera la trayectoria: retirarlo dejó
idénticas 120 corridas de diagnóstico. Por tanto, esas soluciones solo aportan el
punto de partida. Para reproducir el estudio, cada carpeta guarda un manifiesto con
parámetros, semillas, entorno y hashes de código e instancias. Además, un verificador
recalcula cada calendario final con el planificador estricto.

\subsection{Corridas y análisis estadístico}
\label{subsec:estadistica}
El diseño es emparejado: cada método recibe la misma instancia y semilla. El estudio
de congestión usa 48 instancias y 20 semillas, es decir, 960 pares por método. Cada
par ejecuta cuatro trayectorias y reutiliza el prefijo de Tabu uniforme. Por tanto,
las alternativas de un mismo par no son independientes.

La familia confirmatoria tiene cuatro comparaciones. C1 compara D-MAB con Tabu
uniforme y constituye la hipótesis primaria. C2 compara la cartera con Tabu
prolongado, y C3 la compara con dos Tabu. C4 compara D-MAB con UCB1. C1 y C4 usan
3.030 evaluaciones; C2 y C3 usan 6.060. Cada diferencia se calcula como alternativa
menos referencia. Por tanto, un valor negativo favorece a la alternativa.

Las semillas de una instancia comparten datos. Por eso, la prueba de Wilcoxon
bilateral \parencite{wilcoxon1945} opera sobre medias por instancia y no sobre pares.
Holm \parencite{holm1979} corrige las cuatro comparaciones. Un bootstrap jerárquico
\parencite{efron1979} de instancias y semillas, con 5.000 réplicas, entrega
intervalos de 95~\%. El umbral práctico es 0,25 unidades de $f$. Todos los
tiempos son múltiplos de 0,5 minutos, y un paso de espera total mueve $f$ en 0,25.
La hipótesis primaria exige $p$ de Holm menor que 0,05, intervalo bajo cero y media
de a lo sumo $-0{,}25$.

\begin{table}[ht]
 \centering
 \caption{Tamaños para potencia 0,80 y $\alpha=0{,}05$, según diferencias de 240 pares. DE: desviación estándar de la diferencia.}
 \label{tab:potencia}
 \begin{tabular}{@{}lrrrrr@{}}
  \toprule
  & & \multicolumn{2}{c}{Pares necesarios} & \multicolumn{2}{c}{Instancias (20 semillas)}\\
  Comparación & DE & Efecto 0,25 & Efecto 0,65 & Efecto 0,25 & Efecto 0,65\\
  \midrule
  D-MAB $-$ Tabu uniforme & 7,034 & 6.217 & 922 & 317 & 49\\
  Cartera $-$ Tabu prolongado & 4,292 & 2.315 & 345 & 221 & 35\\
  Cartera $-$ Dos Tabu & 5,018 & 3.165 & 470 & 322 & 50\\
  \bottomrule
 \end{tabular}
\end{table}

La Tabla~\ref{tab:potencia} justifica el número de corridas con los datos del estudio
de carteras. El diseño de 240 pares está subdimensionado: C1 necesita 922 pares para
detectar 0,65 y 6.217 para detectar 0,25. Son 3,8 y 25,9 veces los pares disponibles.
Además, con 12 instancias, la potencia simulada de la prueba real es 0,30 y 0,08,
respectivamente.
Por esta razón, una regla pre-registrada fijó 48 instancias. Con ese tamaño, C1 alcanza 0,87
para detectar 0,65. Sin embargo, baja a 0,72 con $\alpha=0{,}0125$, el peor caso de
Holm, y a 0,24 para detectar 0,25. Por tanto, un resultado no significativo no probará
equivalencia. Finalmente, un piloto descriptivo de 40 corridas prueba el protocolo con
igual tiempo de pared, porque los demás estudios igualan evaluaciones.

\subsection{Amenazas a la validez}
\label{subsec:amenazas}
Cinco amenazas limitan la interpretación. Primero, la reparación altera cerca de un
tercio de los candidatos (31,3~\%). Así, la recompensa se atribuye a un movimiento que
no se aplicó tal cual. Segundo, solo 3,9~\% de las propuestas recibe recompensa
positiva, lo que reduce la señal para aprender. Tercero, las instancias estándar ya
eran conocidas al elegir la arquitectura. Cuarto, el catálogo de congestión comparte
el perfil de duraciones del estándar y no constituye un dominio independiente.
Quinto, la potencia simulada supone un efecto homogéneo y extrapola desde 12 instancias.
Además, la DE de C1 en congestión fue 13,21, frente a la DE supuesta de 7,03. Con la
aproximación t, la DE supuesta pedía 49 instancias para potencia 0,80. Con la DE
observada, 48 instancias alcanzan 0,32.

% Usa solo paquetes ya cargados por main.tex y pucv_inf_2024.sty.
% Referencia eq:objetivo, definida en la Propuesta metodológica.
\section{Configuración experimental}
\label{sec:experimental}

Este diseño responde a la pregunta del proyecto: ¿en qué condiciones D-MAB
mejora Tabu Search frente a la selección uniforme? Por ello, las comparaciones
comparten representación, planificador, objetivo y presupuesto. La hipótesis
primaria plantea que, bajo congestión, Tabu con D-MAB obtiene menor $f$ medio que
Tabu uniforme. El equipo fijó hipótesis, semillas y pruebas en un protocolo previo
a los experimentos nuevos. Esta sección actualiza el plan de evaluación de la sección
anterior y prevalece ante cualquier diferencia. Además, describe el diseño y no informa
conclusiones.

\subsection{Conjuntos de datos}
\label{subsec:datos}
El estudio usa dos catálogos sintéticos, sin datos clínicos, y los analiza por
separado. El catálogo estándar proviene del repositorio oficial
\parencite{saicooh2026}. Contiene 12 instancias: tres réplicas de 15, 20, 25 y 30
cirugías. Sin embargo, este catálogo casi no tensiona la espera entre etapas. Solo
16 de 240 soluciones finales del estudio de carteras presentan espera positiva.
Por tanto, no permite evaluar la condición de bloqueo que plantea la pregunta.

Para cubrir ese vacío, el equipo generó un catálogo de congestión de 48 instancias,
doce por tamaño. El generador remuestrea pacientes de una instancia estándar y
aplica ruido lognormal con $\sigma=0{,}10$. Luego redondea los tiempos a 0,5 minutos
y redibuja la elegibilidad de personal. Cada instancia tiene semilla propia.

Una calibración evaluó 18 combinaciones de salas, cirujanos y espera máxima, sobre
cuatro instancias de desarrollo excluidas del catálogo. La regla pre-registrada eligió
reducir los cirujanos al 40~\% y mantener salas y espera máxima. Esta elección
sigue el mecanismo del modelo: con la misma sala en ambas etapas, la espera depende
del cirujano. De hecho, sin reducir cirujanos, ninguna combinación superó 10~\% de
soluciones con espera. En el catálogo final, 51,8~\% de las soluciones de Tabu uniforme
y D-MAB presenta espera positiva. Este dato verifica la presión sobre la restricción.

Ambos catálogos superan la validación de esquema, tiempos, recursos, elegibilidad y
digest canónico. Además, FIFO, SPT y LPT encuentran solución factible en las 48
instancias nuevas. Los manifiestos guardan el SHA-256 de cada archivo. Cada
experimento usa un rango exclusivo de semillas. Las semillas nuevas sobre el
catálogo estándar miden variación de la búsqueda; no miden generalización a
instancias no vistas.

\subsection{Líneas base y controles}
\label{subsec:controles}
La Tabla~\ref{tab:lineas-base} vincula cada control con la explicación alternativa
que descarta. Así, una diferencia observada puede atribuirse al aprendizaje y no a
otro componente del diseño.

\begin{table}[ht]
 \centering
 \caption{Controles y explicación alternativa que descarta cada uno.}
 \label{tab:lineas-base}
 \begin{tabular}{@{}p{0.33\textwidth}p{0.61\textwidth}@{}}
  \toprule
  Control (evaluaciones) & Explicación alternativa\\
  \midrule
  FIFO, SPT y LPT & Una regla constructiva basta. Integran el inicio de Tabu; superarlas no prueba el aporte de D-MAB.\\
  Tabu uniforme (3.030) & La mejora proviene de Tabu y no de aprender.\\
  Tabu prolongado (6.060) & La cartera mejora solo por consumir más evaluaciones.\\
  Dos Tabu independientes (6.060) & La cartera mejora por repetir la búsqueda.\\
  Tabu + UCB1 sin Page--Hinkley (3.030) & El efecto proviene de UCB1 y no del reinicio.\\
  \bottomrule
 \end{tabular}
\end{table}

\subsection{Métricas}
\label{subsec:metricas}
La métrica primaria es $f(S)$ de la ecuación~\eqref{eq:objetivo}, fijada antes de
evaluar. Las métricas secundarias son $C_{\max}$, espera total, espera máxima,
factibilidad, tiempo de pared y convergencia. La convergencia se mide con la
evaluación de la última mejora global. Estas métricas son descriptivas. Una baja de
$f$ no implica una baja de cada componente; por eso, se informan por separado.

\subsection{Implementación}
\label{subsec:implementacion}
El código usa Python~3.13.2, NumPy~2.2.3, SciPy~1.15.2 y PyYAML~6.0.3. Los
experimentos nuevos corren en un AMD Ryzen~7 8845HS, con 8 núcleos, 16 hilos y
13,8~GiB de RAM, sobre Windows~11. La Tabla~\ref{tab:parametros} fija los parámetros.
El tiempo se mide en un solo proceso. Una cartera tarda, en mediana, 4,12~s con 15
cirugías (rango 3,83--4,80) y 13,09~s con 30 (rango 12,65--13,53). Sin embargo, el
rendimiento varió entre sesiones: con 15 cirugías pasó de 1.470 a cerca de 2.240
evaluaciones por segundo. Por eso, el control principal de cómputo son las evaluaciones.

\begin{table}[ht]
 \centering
 \caption{Parámetros congelados de la búsqueda y del presupuesto.}
 \label{tab:parametros}
 \begin{tabular}{@{}p{0.16\textwidth}p{0.78\textwidth}@{}}
  \toprule
  Componente & Configuración\\
  \midrule
  Tabu & 30 soluciones iniciales (FIFO, SPT, LPT y 27 aleatorias); lotes de 15 candidatos; memoria de 7 firmas de solución; aspiración por mejora global.\\
  D-MAB & Brazos \texttt{swap}, \texttt{insert}, \texttt{room\_1}, \texttt{room\_2} y \texttt{room\_pair}; $C=0{,}1$; $\delta=0{,}01$; $\lambda=1$.\\
  Cartera & Dos caminos de 3.030 evaluaciones; reparto 50/50 fijo; conserva el menor $f$.\\
  Parada & Agotar el presupuesto de evaluaciones del camino.\\
  Conteo & Cuentan las 30 iniciales y cada candidato distinto del lote, aun infactible. No cuentan duplicados, reparación ni revalidación final.\\
  \bottomrule
 \end{tabular}
\end{table}

El archivo de 30 soluciones del código no altera la trayectoria: retirarlo dejó
idénticas 120 corridas de diagnóstico. Por tanto, esas soluciones solo aportan el
punto de partida. Para reproducir el estudio, cada carpeta guarda un manifiesto con
parámetros, semillas, entorno y hashes de código e instancias. Además, un verificador
recalcula cada calendario final con el planificador estricto.

\subsection{Corridas y análisis estadístico}
\label{subsec:estadistica}
El diseño es emparejado: cada método recibe la misma instancia y semilla. El estudio
de congestión usa 48 instancias y 20 semillas, es decir, 960 pares por método. Cada
par ejecuta cuatro trayectorias y reutiliza el prefijo de Tabu uniforme. Por tanto,
las alternativas de un mismo par no son independientes.

La familia confirmatoria tiene cuatro comparaciones. C1 compara D-MAB con Tabu
uniforme y constituye la hipótesis primaria. C2 compara la cartera con Tabu
prolongado, y C3 la compara con dos Tabu. C4 compara D-MAB con UCB1. C1 y C4 usan
3.030 evaluaciones; C2 y C3 usan 6.060. Cada diferencia se calcula como alternativa
menos referencia. Por tanto, un valor negativo favorece a la alternativa.

Las semillas de una instancia comparten datos. Por eso, la prueba de Wilcoxon
bilateral \parencite{wilcoxon1945} opera sobre medias por instancia y no sobre pares.
Holm \parencite{holm1979} corrige las cuatro comparaciones. Un bootstrap jerárquico
\parencite{efron1979} de instancias y semillas, con 5.000 réplicas, entrega
intervalos de 95~\%. El umbral práctico es 0,25 unidades de $f$. Todos los
tiempos son múltiplos de 0,5 minutos, y un paso de espera total mueve $f$ en 0,25.
La hipótesis primaria exige $p$ de Holm menor que 0,05, intervalo bajo cero y media
de a lo sumo $-0{,}25$.

\begin{table}[ht]
 \centering
 \caption{Tamaños para potencia 0,80 y $\alpha=0{,}05$, según diferencias de 240 pares. DE: desviación estándar de la diferencia.}
 \label{tab:potencia}
 \begin{tabular}{@{}lrrrrr@{}}
  \toprule
  & & \multicolumn{2}{c}{Pares necesarios} & \multicolumn{2}{c}{Instancias (20 semillas)}\\
  Comparación & DE & Efecto 0,25 & Efecto 0,65 & Efecto 0,25 & Efecto 0,65\\
  \midrule
  D-MAB $-$ Tabu uniforme & 7,034 & 6.217 & 922 & 317 & 49\\
  Cartera $-$ Tabu prolongado & 4,292 & 2.315 & 345 & 221 & 35\\
  Cartera $-$ Dos Tabu & 5,018 & 3.165 & 470 & 322 & 50\\
  \bottomrule
 \end{tabular}
\end{table}

La Tabla~\ref{tab:potencia} justifica el número de corridas con los datos del estudio
de carteras. El diseño de 240 pares está subdimensionado: C1 necesita 922 pares para
detectar 0,65 y 6.217 para detectar 0,25. Son 3,8 y 25,9 veces los pares disponibles.
Además, con 12 instancias, la potencia simulada de la prueba real es 0,30 y 0,08,
respectivamente.
Por esta razón, una regla pre-registrada fijó 48 instancias. Con ese tamaño, C1 alcanza 0,87
para detectar 0,65. Sin embargo, baja a 0,72 con $\alpha=0{,}0125$, el peor caso de
Holm, y a 0,24 para detectar 0,25. Por tanto, un resultado no significativo no probará
equivalencia. Finalmente, un piloto descriptivo de 40 corridas prueba el protocolo con
igual tiempo de pared, porque los demás estudios igualan evaluaciones.

\subsection{Amenazas a la validez}
\label{subsec:amenazas}
Cinco amenazas limitan la interpretación. Primero, la reparación altera cerca de un
tercio de los candidatos (31,3~\%). Así, la recompensa se atribuye a un movimiento que
no se aplicó tal cual. Segundo, solo 3,9~\% de las propuestas recibe recompensa
positiva, lo que reduce la señal para aprender. Tercero, las instancias estándar ya
eran conocidas al elegir la arquitectura. Cuarto, el catálogo de congestión comparte
el perfil de duraciones del estándar y no constituye un dominio independiente.
Quinto, la potencia simulada supone un efecto homogéneo y extrapola desde 12 instancias.
Además, la DE de C1 en congestión fue 13,21, frente a la DE supuesta de 7,03. Con la
aproximación t, la DE supuesta pedía 49 instancias para potencia 0,80. Con la DE
observada, 48 instancias alcanzan 0,32.

% Usa solo paquetes ya cargados por main.tex y pucv_inf_2024.sty.
% Referencia eq:objetivo, definida en la Propuesta metodológica.
\section{Configuración experimental}
\label{sec:experimental}

Este diseño responde a la pregunta del proyecto: ¿en qué condiciones D-MAB
mejora Tabu Search frente a la selección uniforme? Por ello, las comparaciones
comparten representación, planificador, objetivo y presupuesto. La hipótesis
primaria plantea que, bajo congestión, Tabu con D-MAB obtiene menor $f$ medio que
Tabu uniforme. El equipo fijó hipótesis, semillas y pruebas en un protocolo previo
a los experimentos nuevos. Esta sección actualiza el plan de evaluación de la sección
anterior y prevalece ante cualquier diferencia. Además, describe el diseño y no informa
conclusiones.

\subsection{Conjuntos de datos}
\label{subsec:datos}
El estudio usa dos catálogos sintéticos, sin datos clínicos, y los analiza por
separado. El catálogo estándar proviene del repositorio oficial
\parencite{saicooh2026}. Contiene 12 instancias: tres réplicas de 15, 20, 25 y 30
cirugías. Sin embargo, este catálogo casi no tensiona la espera entre etapas. Solo
16 de 240 soluciones finales del estudio de carteras presentan espera positiva.
Por tanto, no permite evaluar la condición de bloqueo que plantea la pregunta.

Para cubrir ese vacío, el equipo generó un catálogo de congestión de 48 instancias,
doce por tamaño. El generador remuestrea pacientes de una instancia estándar y
aplica ruido lognormal con $\sigma=0{,}10$. Luego redondea los tiempos a 0,5 minutos
y redibuja la elegibilidad de personal. Cada instancia tiene semilla propia.

Una calibración evaluó 18 combinaciones de salas, cirujanos y espera máxima, sobre
cuatro instancias de desarrollo excluidas del catálogo. La regla pre-registrada eligió
reducir los cirujanos al 40~\% y mantener salas y espera máxima. Esta elección
sigue el mecanismo del modelo: con la misma sala en ambas etapas, la espera depende
del cirujano. De hecho, sin reducir cirujanos, ninguna combinación superó 10~\% de
soluciones con espera. En el catálogo final, 51,8~\% de las soluciones de Tabu uniforme
y D-MAB presenta espera positiva. Este dato verifica la presión sobre la restricción.

Ambos catálogos superan la validación de esquema, tiempos, recursos, elegibilidad y
digest canónico. Además, FIFO, SPT y LPT encuentran solución factible en las 48
instancias nuevas. Los manifiestos guardan el SHA-256 de cada archivo. Cada
experimento usa un rango exclusivo de semillas. Las semillas nuevas sobre el
catálogo estándar miden variación de la búsqueda; no miden generalización a
instancias no vistas.

\subsection{Líneas base y controles}
\label{subsec:controles}
La Tabla~\ref{tab:lineas-base} vincula cada control con la explicación alternativa
que descarta. Así, una diferencia observada puede atribuirse al aprendizaje y no a
otro componente del diseño.

\begin{table}[ht]
 \centering
 \caption{Controles y explicación alternativa que descarta cada uno.}
 \label{tab:lineas-base}
 \begin{tabular}{@{}p{0.33\textwidth}p{0.61\textwidth}@{}}
  \toprule
  Control (evaluaciones) & Explicación alternativa\\
  \midrule
  FIFO, SPT y LPT & Una regla constructiva basta. Integran el inicio de Tabu; superarlas no prueba el aporte de D-MAB.\\
  Tabu uniforme (3.030) & La mejora proviene de Tabu y no de aprender.\\
  Tabu prolongado (6.060) & La cartera mejora solo por consumir más evaluaciones.\\
  Dos Tabu independientes (6.060) & La cartera mejora por repetir la búsqueda.\\
  Tabu + UCB1 sin Page--Hinkley (3.030) & El efecto proviene de UCB1 y no del reinicio.\\
  \bottomrule
 \end{tabular}
\end{table}

\subsection{Métricas}
\label{subsec:metricas}
La métrica primaria es $f(S)$ de la ecuación~\eqref{eq:objetivo}, fijada antes de
evaluar. Las métricas secundarias son $C_{\max}$, espera total, espera máxima,
factibilidad, tiempo de pared y convergencia. La convergencia se mide con la
evaluación de la última mejora global. Estas métricas son descriptivas. Una baja de
$f$ no implica una baja de cada componente; por eso, se informan por separado.

\subsection{Implementación}
\label{subsec:implementacion}
El código usa Python~3.13.2, NumPy~2.2.3, SciPy~1.15.2 y PyYAML~6.0.3. Los
experimentos nuevos corren en un AMD Ryzen~7 8845HS, con 8 núcleos, 16 hilos y
13,8~GiB de RAM, sobre Windows~11. La Tabla~\ref{tab:parametros} fija los parámetros.
El tiempo se mide en un solo proceso. Una cartera tarda, en mediana, 4,12~s con 15
cirugías (rango 3,83--4,80) y 13,09~s con 30 (rango 12,65--13,53). Sin embargo, el
rendimiento varió entre sesiones: con 15 cirugías pasó de 1.470 a cerca de 2.240
evaluaciones por segundo. Por eso, el control principal de cómputo son las evaluaciones.

\begin{table}[ht]
 \centering
 \caption{Parámetros congelados de la búsqueda y del presupuesto.}
 \label{tab:parametros}
 \begin{tabular}{@{}p{0.16\textwidth}p{0.78\textwidth}@{}}
  \toprule
  Componente & Configuración\\
  \midrule
  Tabu & 30 soluciones iniciales (FIFO, SPT, LPT y 27 aleatorias); lotes de 15 candidatos; memoria de 7 firmas de solución; aspiración por mejora global.\\
  D-MAB & Brazos \texttt{swap}, \texttt{insert}, \texttt{room\_1}, \texttt{room\_2} y \texttt{room\_pair}; $C=0{,}1$; $\delta=0{,}01$; $\lambda=1$.\\
  Cartera & Dos caminos de 3.030 evaluaciones; reparto 50/50 fijo; conserva el menor $f$.\\
  Parada & Agotar el presupuesto de evaluaciones del camino.\\
  Conteo & Cuentan las 30 iniciales y cada candidato distinto del lote, aun infactible. No cuentan duplicados, reparación ni revalidación final.\\
  \bottomrule
 \end{tabular}
\end{table}

El archivo de 30 soluciones del código no altera la trayectoria: retirarlo dejó
idénticas 120 corridas de diagnóstico. Por tanto, esas soluciones solo aportan el
punto de partida. Para reproducir el estudio, cada carpeta guarda un manifiesto con
parámetros, semillas, entorno y hashes de código e instancias. Además, un verificador
recalcula cada calendario final con el planificador estricto.

\subsection{Corridas y análisis estadístico}
\label{subsec:estadistica}
El diseño es emparejado: cada método recibe la misma instancia y semilla. El estudio
de congestión usa 48 instancias y 20 semillas, es decir, 960 pares por método. Cada
par ejecuta cuatro trayectorias y reutiliza el prefijo de Tabu uniforme. Por tanto,
las alternativas de un mismo par no son independientes.

La familia confirmatoria tiene cuatro comparaciones. C1 compara D-MAB con Tabu
uniforme y constituye la hipótesis primaria. C2 compara la cartera con Tabu
prolongado, y C3 la compara con dos Tabu. C4 compara D-MAB con UCB1. C1 y C4 usan
3.030 evaluaciones; C2 y C3 usan 6.060. Cada diferencia se calcula como alternativa
menos referencia. Por tanto, un valor negativo favorece a la alternativa.

Las semillas de una instancia comparten datos. Por eso, la prueba de Wilcoxon
bilateral \parencite{wilcoxon1945} opera sobre medias por instancia y no sobre pares.
Holm \parencite{holm1979} corrige las cuatro comparaciones. Un bootstrap jerárquico
\parencite{efron1979} de instancias y semillas, con 5.000 réplicas, entrega
intervalos de 95~\%. El umbral práctico es 0,25 unidades de $f$. Todos los
tiempos son múltiplos de 0,5 minutos, y un paso de espera total mueve $f$ en 0,25.
La hipótesis primaria exige $p$ de Holm menor que 0,05, intervalo bajo cero y media
de a lo sumo $-0{,}25$.

\begin{table}[ht]
 \centering
 \caption{Tamaños para potencia 0,80 y $\alpha=0{,}05$, según diferencias de 240 pares. DE: desviación estándar de la diferencia.}
 \label{tab:potencia}
 \begin{tabular}{@{}lrrrrr@{}}
  \toprule
  & & \multicolumn{2}{c}{Pares necesarios} & \multicolumn{2}{c}{Instancias (20 semillas)}\\
  Comparación & DE & Efecto 0,25 & Efecto 0,65 & Efecto 0,25 & Efecto 0,65\\
  \midrule
  D-MAB $-$ Tabu uniforme & 7,034 & 6.217 & 922 & 317 & 49\\
  Cartera $-$ Tabu prolongado & 4,292 & 2.315 & 345 & 221 & 35\\
  Cartera $-$ Dos Tabu & 5,018 & 3.165 & 470 & 322 & 50\\
  \bottomrule
 \end{tabular}
\end{table}

La Tabla~\ref{tab:potencia} justifica el número de corridas con los datos del estudio
de carteras. El diseño de 240 pares está subdimensionado: C1 necesita 922 pares para
detectar 0,65 y 6.217 para detectar 0,25. Son 3,8 y 25,9 veces los pares disponibles.
Además, con 12 instancias, la potencia simulada de la prueba real es 0,30 y 0,08,
respectivamente.
Por esta razón, una regla pre-registrada fijó 48 instancias. Con ese tamaño, C1 alcanza 0,87
para detectar 0,65. Sin embargo, baja a 0,72 con $\alpha=0{,}0125$, el peor caso de
Holm, y a 0,24 para detectar 0,25. Por tanto, un resultado no significativo no probará
equivalencia. Finalmente, un piloto descriptivo de 40 corridas prueba el protocolo con
igual tiempo de pared, porque los demás estudios igualan evaluaciones.

\subsection{Amenazas a la validez}
\label{subsec:amenazas}
Cinco amenazas limitan la interpretación. Primero, la reparación altera cerca de un
tercio de los candidatos (31,3~\%). Así, la recompensa se atribuye a un movimiento que
no se aplicó tal cual. Segundo, solo 3,9~\% de las propuestas recibe recompensa
positiva, lo que reduce la señal para aprender. Tercero, las instancias estándar ya
eran conocidas al elegir la arquitectura. Cuarto, el catálogo de congestión comparte
el perfil de duraciones del estándar y no constituye un dominio independiente.
Quinto, la potencia simulada supone un efecto homogéneo y extrapola desde 12 instancias.
Además, la DE de C1 en congestión fue 13,21, frente a la DE supuesta de 7,03. Con la
aproximación t, la DE supuesta pedía 49 instancias para potencia 0,80. Con la DE
observada, 48 instancias alcanzan 0,32.

% Usa solo paquetes ya cargados por main.tex y pucv_inf_2024.sty.
% Referencia eq:objetivo, definida en la Propuesta metodológica.
\section{Configuración experimental}
\label{sec:experimental}

Este diseño responde a la pregunta del proyecto: ¿en qué condiciones D-MAB
mejora Tabu Search frente a la selección uniforme? Por ello, las comparaciones
comparten representación, planificador, objetivo y presupuesto. La hipótesis
primaria plantea que, bajo congestión, Tabu con D-MAB obtiene menor $f$ medio que
Tabu uniforme. El equipo fijó hipótesis, semillas y pruebas en un protocolo previo
a los experimentos nuevos. Esta sección actualiza el plan de evaluación de la sección
anterior y prevalece ante cualquier diferencia. Además, describe el diseño y no informa
conclusiones.

\subsection{Conjuntos de datos}
\label{subsec:datos}
El estudio usa dos catálogos sintéticos, sin datos clínicos, y los analiza por
separado. El catálogo estándar proviene del repositorio oficial
\parencite{saicooh2026}. Contiene 12 instancias: tres réplicas de 15, 20, 25 y 30
cirugías. Sin embargo, este catálogo casi no tensiona la espera entre etapas. Solo
16 de 240 soluciones finales del estudio de carteras presentan espera positiva.
Por tanto, no permite evaluar la condición de bloqueo que plantea la pregunta.

Para cubrir ese vacío, el equipo generó un catálogo de congestión de 48 instancias,
doce por tamaño. El generador remuestrea pacientes de una instancia estándar y
aplica ruido lognormal con $\sigma=0{,}10$. Luego redondea los tiempos a 0,5 minutos
y redibuja la elegibilidad de personal. Cada instancia tiene semilla propia.

Una calibración evaluó 18 combinaciones de salas, cirujanos y espera máxima, sobre
cuatro instancias de desarrollo excluidas del catálogo. La regla pre-registrada eligió
reducir los cirujanos al 40~\% y mantener salas y espera máxima. Esta elección
sigue el mecanismo del modelo: con la misma sala en ambas etapas, la espera depende
del cirujano. De hecho, sin reducir cirujanos, ninguna combinación superó 10~\% de
soluciones con espera. En el catálogo final, 51,8~\% de las soluciones de Tabu uniforme
y D-MAB presenta espera positiva. Este dato verifica la presión sobre la restricción.

Ambos catálogos superan la validación de esquema, tiempos, recursos, elegibilidad y
digest canónico. Además, FIFO, SPT y LPT encuentran solución factible en las 48
instancias nuevas. Los manifiestos guardan el SHA-256 de cada archivo. Cada
experimento usa un rango exclusivo de semillas. Las semillas nuevas sobre el
catálogo estándar miden variación de la búsqueda; no miden generalización a
instancias no vistas.

\subsection{Líneas base y controles}
\label{subsec:controles}
La Tabla~\ref{tab:lineas-base} vincula cada control con la explicación alternativa
que descarta. Así, una diferencia observada puede atribuirse al aprendizaje y no a
otro componente del diseño.

\begin{table}[ht]
 \centering
 \caption{Controles y explicación alternativa que descarta cada uno.}
 \label{tab:lineas-base}
 \begin{tabular}{@{}p{0.33\textwidth}p{0.61\textwidth}@{}}
  \toprule
  Control (evaluaciones) & Explicación alternativa\\
  \midrule
  FIFO, SPT y LPT & Una regla constructiva basta. Integran el inicio de Tabu; superarlas no prueba el aporte de D-MAB.\\
  Tabu uniforme (3.030) & La mejora proviene de Tabu y no de aprender.\\
  Tabu prolongado (6.060) & La cartera mejora solo por consumir más evaluaciones.\\
  Dos Tabu independientes (6.060) & La cartera mejora por repetir la búsqueda.\\
  Tabu + UCB1 sin Page--Hinkley (3.030) & El efecto proviene de UCB1 y no del reinicio.\\
  \bottomrule
 \end{tabular}
\end{table}

\subsection{Métricas}
\label{subsec:metricas}
La métrica primaria es $f(S)$ de la ecuación~\eqref{eq:objetivo}, fijada antes de
evaluar. Las métricas secundarias son $C_{\max}$, espera total, espera máxima,
factibilidad, tiempo de pared y convergencia. La convergencia se mide con la
evaluación de la última mejora global. Estas métricas son descriptivas. Una baja de
$f$ no implica una baja de cada componente; por eso, se informan por separado.

\subsection{Implementación}
\label{subsec:implementacion}
El código usa Python~3.13.2, NumPy~2.2.3, SciPy~1.15.2 y PyYAML~6.0.3. Los
experimentos nuevos corren en un AMD Ryzen~7 8845HS, con 8 núcleos, 16 hilos y
13,8~GiB de RAM, sobre Windows~11. La Tabla~\ref{tab:parametros} fija los parámetros.
El tiempo se mide en un solo proceso. Una cartera tarda, en mediana, 4,12~s con 15
cirugías (rango 3,83--4,80) y 13,09~s con 30 (rango 12,65--13,53). Sin embargo, el
rendimiento varió entre sesiones: con 15 cirugías pasó de 1.470 a cerca de 2.240
evaluaciones por segundo. Por eso, el control principal de cómputo son las evaluaciones.

\begin{table}[ht]
 \centering
 \caption{Parámetros congelados de la búsqueda y del presupuesto.}
 \label{tab:parametros}
 \begin{tabular}{@{}p{0.16\textwidth}p{0.78\textwidth}@{}}
  \toprule
  Componente & Configuración\\
  \midrule
  Tabu & 30 soluciones iniciales (FIFO, SPT, LPT y 27 aleatorias); lotes de 15 candidatos; memoria de 7 firmas de solución; aspiración por mejora global.\\
  D-MAB & Brazos \texttt{swap}, \texttt{insert}, \texttt{room\_1}, \texttt{room\_2} y \texttt{room\_pair}; $C=0{,}1$; $\delta=0{,}01$; $\lambda=1$.\\
  Cartera & Dos caminos de 3.030 evaluaciones; reparto 50/50 fijo; conserva el menor $f$.\\
  Parada & Agotar el presupuesto de evaluaciones del camino.\\
  Conteo & Cuentan las 30 iniciales y cada candidato distinto del lote, aun infactible. No cuentan duplicados, reparación ni revalidación final.\\
  \bottomrule
 \end{tabular}
\end{table}

El archivo de 30 soluciones del código no altera la trayectoria: retirarlo dejó
idénticas 120 corridas de diagnóstico. Por tanto, esas soluciones solo aportan el
punto de partida. Para reproducir el estudio, cada carpeta guarda un manifiesto con
parámetros, semillas, entorno y hashes de código e instancias. Además, un verificador
recalcula cada calendario final con el planificador estricto.

\subsection{Corridas y análisis estadístico}
\label{subsec:estadistica}
El diseño es emparejado: cada método recibe la misma instancia y semilla. El estudio
de congestión usa 48 instancias y 20 semillas, es decir, 960 pares por método. Cada
par ejecuta cuatro trayectorias y reutiliza el prefijo de Tabu uniforme. Por tanto,
las alternativas de un mismo par no son independientes.

La familia confirmatoria tiene cuatro comparaciones. C1 compara D-MAB con Tabu
uniforme y constituye la hipótesis primaria. C2 compara la cartera con Tabu
prolongado, y C3 la compara con dos Tabu. C4 compara D-MAB con UCB1. C1 y C4 usan
3.030 evaluaciones; C2 y C3 usan 6.060. Cada diferencia se calcula como alternativa
menos referencia. Por tanto, un valor negativo favorece a la alternativa.

Las semillas de una instancia comparten datos. Por eso, la prueba de Wilcoxon
bilateral \parencite{wilcoxon1945} opera sobre medias por instancia y no sobre pares.
Holm \parencite{holm1979} corrige las cuatro comparaciones. Un bootstrap jerárquico
\parencite{efron1979} de instancias y semillas, con 5.000 réplicas, entrega
intervalos de 95~\%. El umbral práctico es 0,25 unidades de $f$. Todos los
tiempos son múltiplos de 0,5 minutos, y un paso de espera total mueve $f$ en 0,25.
La hipótesis primaria exige $p$ de Holm menor que 0,05, intervalo bajo cero y media
de a lo sumo $-0{,}25$.

\begin{table}[ht]
 \centering
 \caption{Tamaños para potencia 0,80 y $\alpha=0{,}05$, según diferencias de 240 pares. DE: desviación estándar de la diferencia.}
 \label{tab:potencia}
 \begin{tabular}{@{}lrrrrr@{}}
  \toprule
  & & \multicolumn{2}{c}{Pares necesarios} & \multicolumn{2}{c}{Instancias (20 semillas)}\\
  Comparación & DE & Efecto 0,25 & Efecto 0,65 & Efecto 0,25 & Efecto 0,65\\
  \midrule
  D-MAB $-$ Tabu uniforme & 7,034 & 6.217 & 922 & 317 & 49\\
  Cartera $-$ Tabu prolongado & 4,292 & 2.315 & 345 & 221 & 35\\
  Cartera $-$ Dos Tabu & 5,018 & 3.165 & 470 & 322 & 50\\
  \bottomrule
 \end{tabular}
\end{table}

La Tabla~\ref{tab:potencia} justifica el número de corridas con los datos del estudio
de carteras. El diseño de 240 pares está subdimensionado: C1 necesita 922 pares para
detectar 0,65 y 6.217 para detectar 0,25. Son 3,8 y 25,9 veces los pares disponibles.
Además, con 12 instancias, la potencia simulada de la prueba real es 0,30 y 0,08,
respectivamente.
Por esta razón, una regla pre-registrada fijó 48 instancias. Con ese tamaño, C1 alcanza 0,87
para detectar 0,65. Sin embargo, baja a 0,72 con $\alpha=0{,}0125$, el peor caso de
Holm, y a 0,24 para detectar 0,25. Por tanto, un resultado no significativo no probará
equivalencia. Finalmente, un piloto descriptivo de 40 corridas prueba el protocolo con
igual tiempo de pared, porque los demás estudios igualan evaluaciones.

\subsection{Amenazas a la validez}
\label{subsec:amenazas}
Cinco amenazas limitan la interpretación. Primero, la reparación altera cerca de un
tercio de los candidatos (31,3~\%). Así, la recompensa se atribuye a un movimiento que
no se aplicó tal cual. Segundo, solo 3,9~\% de las propuestas recibe recompensa
positiva, lo que reduce la señal para aprender. Tercero, las instancias estándar ya
eran conocidas al elegir la arquitectura. Cuarto, el catálogo de congestión comparte
el perfil de duraciones del estándar y no constituye un dominio independiente.
Quinto, la potencia simulada supone un efecto homogéneo y extrapola desde 12 instancias.
Además, la DE de C1 en congestión fue 13,21, frente a la DE supuesta de 7,03. Con la
aproximación t, la DE supuesta pedía 49 instancias para potencia 0,80. Con la DE
observada, 48 instancias alcanzan 0,32.
